\documentclass[12pt,a4paper]{article}
\usepackage[utf8]{inputenc}
\usepackage[T1]{fontenc}
\usepackage[turkish]{babel}
\usepackage{geometry}
\usepackage{booktabs}
\usepackage{array}
\usepackage{longtable}
\usepackage{hyperref}
\usepackage{setspace}
\geometry{margin=2.5cm}
\setstretch{1.15}

\title{Haber Ba\c{s}l\i{}klar\i{}nda \c{C}o\u{g}altma Tespiti:\\
Sistem Raporu ve Zipf Analizi}
\author{Mustafa Mete (2202131015) \and Emine Avcu (2202131006)}
\date{\today}

\begin{document}
\maketitle

\section{Proje \.Ozeti}
Bu \c{c}al\i{}\c{s}mada farkl\i{} kaynaklardan \.{I}ngilizce haber ba\c{s}l\i{}klar\i{} toplanm\i{}\c{s}, tek bir veri havuzunda birle\c{s}tirilmi\c{s}, normalizasyon uygulanm\i{}\c{s} ve benzer (muhtemel duplicate) ba\c{s}l\i{}klar TF-IDF + cosine similarity ile tespit edilmi\c{s}tir. Ek olarak veri kalitesi ve dilsel da\u{g}\i{}l\i{}m i\c{c}in Zipf analizi ger\c{c}ekle\c{s}tirilmi\c{s}tir.

\section{Repo Denetimi (\.Inceleme Kapsam\i{})}
Bu rapor haz\i{}rlan\i{}rken depo i\c{c}eri\u{g}i dosya baz\i{}nda taranm\i{}\c{s}, mevcut kod/konfig\"urasyon ve \"uretilen \c{c}\i{}kt\i{}lar\i{}n son durumu okunarak \c{c}\i{}kar\i{}m yap\i{}lm\i{}\c{s}t\i{}r.

\subsection{Kimler Haz\i{}rlad\i{}?}
Bu projeyi haz\i{}rlayan ekip:
\begin{itemize}
    \item Mustafa Mete (2202131015)
    \item Emine Avcu (2202131006)
\end{itemize}

\subsection{Dosya Envanteri (Kod ve Ana Dosyalar)}
\begin{itemize}
    \item Python fetch botu say\i{}s\i{}: \textbf{3}
    \item Notebook say\i{}s\i{} (ana): \textbf{1} (\texttt{analysis\_pipeline.ipynb})
    \item Ana konfig\"urasyon: \texttt{.env}
    \item Dok\"umantasyon: \texttt{README.md}
    \item Ba\u{g}\i{}ml\i{}l\i{}klar: \texttt{requirements.txt}
\end{itemize}


\section{Sistem Mimarisi}
\subsection{Veri Toplama Katman\i{}}
Veri toplama botlar\i{} yaln\i{}zca \texttt{bots/fetch} alt\i{}nda konumland\i{}r\i{}lm\i{}\c{s}t\i{}r:
\begin{itemize}
    \item \texttt{newsdata\_fetch\_bot.py}
    \item \texttt{newsapi\_fetch\_bot.py}
    \item \texttt{topic\_search\_fetch\_bot.py}
\end{itemize}

T\"um botlar \.{I}ngilizce toplama kural\i{} ile (\texttt{en}/\texttt{US}) \c{c}al\i{}\c{s}acak \c{s}ekilde kurgulanm\i{}\c{s} ve veriyi ayn\i{} hedef dosyalara eklemeli (append) yazar:
\begin{itemize}
    \item \texttt{data/all/headlines\_all.json}
    \item \texttt{data/all/headlines\_all.csv}
\end{itemize}

\subsection{Analiz Katman\i{}}
Temizleme, normalizasyon, duplicate tespiti ve Zipf analizi tek notebook i\c{c}inde toplanm\i{}\c{s}t\i{}r:
\begin{itemize}
    \item \texttt{analysis\_pipeline.ipynb}
\end{itemize}
Bu notebook, ham veriyi okuyup i\c{s}leyerek a\c{s}a\u{g}\i{}daki \c{c}\i{}kt\i{}lar\i{} \"uretir:
\begin{itemize}
    \item \texttt{data/processed/all/headlines\_normalized\_all.json}
    \item \texttt{data/processed/all/headlines\_normalized\_all.csv}
    \item \texttt{data/results/all/duplicate\_pairs\_all.json}
    \item \texttt{data/results/all/duplicate\_pairs\_all.csv}
\end{itemize}

\section{Veri ve Say\i{}sal Bulgular}
A\c{s}a\u{g}\i{}daki de\u{g}erler mevcut \c{c}\i{}kt\i{} dosyalar\i{}ndan al\i{}nm\i{}\c{s}t\i{}r:
\begin{itemize}
    \item Ham kay\i{}t say\i{}s\i{} (\texttt{headlines\_all.json}): \textbf{2805}
    \item Normalize kay\i{}t say\i{}s\i{} (\texttt{headlines\_normalized\_all.json}): \textbf{2798}
    \item Duplicate \c{c}ift say\i{}s\i{} (\texttt{duplicate\_pairs\_all.json}): \textbf{161}
    \item Toplam token say\i{}s\i{}: \textbf{21428}
    \item Tekil kelime (vocabulary) boyutu: \textbf{6619}
\end{itemize}

\subsection{Kaynak Da\u{g}\i{}l\i{}m\i{} (ilk 10)}
\begin{center}
\begin{tabular}{lr}
\toprule
Kaynak & Kay\i{}t Say\i{}s\i{} \\
\midrule
Yahoo Finance & 96 \\
Let's Data Science & 81 \\
Nature & 60 \\
Business Insider & 44 \\
The Motley Fool & 42 \\
Google News & 33 \\
BBC & 28 \\
The New York Times & 27 \\
Fortune & 25 \\
CNBC & 24 \\
\bottomrule
\end{tabular}
\end{center}

\noindent
Not: \.{I}lk 10 d\i{}\c{s}\i{}nda da Reuters (24), PR Newswire (21), Benzinga (21), MarketBeat (19), The Guardian (18) gibi kaynaklar veri havuzuna katk\i{} sa\u{g}lam\i{}\c{s}t\i{}r.

\subsection{En S\i{}k Kelimeler (ilk 15)}
\begin{center}
\begin{tabular}{lr}
\toprule
Kelime & Frekans \\
\midrule
technology & 153 \\
stock & 147 \\
market & 131 \\
global & 118 \\
innovation & 111 \\
energy & 107 \\
quantum & 107 \\
space & 104 \\
computing & 103 \\
cloud & 95 \\
business & 94 \\
iran & 93 \\
world & 88 \\
climate & 87 \\
software & 87 \\
\bottomrule
\end{tabular}
\end{center}

\subsection{Duplicate Benzerlik \.{I}statistikleri}
Duplicate olarak i\c{s}aretlenen \c{c}iftler i\c{c}in similarity de\u{g}erleri:
\begin{itemize}
    \item Minimum: \textbf{0.8608}
    \item Maksimum: \textbf{1.0000}
    \item Ortalama: \textbf{0.9846}
    \item Medyan: \textbf{1.0000}
\end{itemize}
Bu sonu\c{c}lar, se\c{c}ilen e\c{s}ik (\texttt{0.85}) sonras\i{}nda olduk\c{c}a y\"uksek benzerlikli \c{c}iftlerin kald\i{}\u{g}\i{}n\i{} g\"ostermektedir.

\subsection{Sonu\c{c} Dosyas\i{}ndan \"Ornek \c{C}iftler ve Nedenleri}
\texttt{data/results/all/duplicate\_pairs\_all.json} dosyas\i{}ndan \"ornekler:

\begin{enumerate}
    \item \textbf{``In Guatemala, new AI technology will be listening for illegal deforestation''} \\
    Similarity: \textbf{0.9181} \\
    \textit{Neden:} Ayn\i{} i\c{c}erik birinde orijinal kaynak ba\c{s}l\i{}\u{g}\i{}, di\u{g}erinde Google News/RSS \"uzerinden ``- news - Mongabay'' ekiyle ge\c{c}iyor. Normalizasyon ekleri temizledi\u{g}i i\c{c}in metinler yak\i{}nla\c{s}\i{}yor.

    \item \textbf{``Fairfax India Holdings Corporation: First Quarter Financial Results''} \\
    Similarity: \textbf{1.0000} \\
    \textit{Neden:} Ayn\i{} ba\c{s}l\i{}k farkl\i{} haber da\u{g}\i{}t\i{}m kanallar\i{}nda (Benzinga ve GlobeNewswire) birebir yay\i{}nlanm\i{}\c{s}. Klasik sendikasyon/press release kopyas\i{}.

    \item \textbf{``Fed chair Jerome Powell says ...''} (Google News vs ABC News) \\
    Similarity: \textbf{0.9309} \\
    \textit{Neden:} Ayn\i{} haber c\"umlesine bir kaynakta ``Breaking News, Latest News and Videos'' gibi \c{c}evresel ifade eklenmi\c{s}. Normalizasyon sonras\i{} omurga metin ayn\i{} kald\i{}\u{g}\i{} i\c{c}in duplicate olarak i\c{s}aretleniyor.

    \item \textbf{``News | ...''} \textit{vs} \textbf{``E-News | ...''} \\
    Similarity: \textbf{1.0000} (\texttt{normalized\_title = "news"}) \\
    \textit{Neden:} Bu bir \textbf{yanl\i{}\c{s} pozitif} \"orne\u{g}idir. A\c{s}\i{}r\i{} agresif temizleme sonucunda farkl\i{} ba\c{s}l\i{}klar ``news'' tokenina kadar d\"u\c{s}m\"u\c{s}. Bu nedenle alakas\i{}z i\c{c}erikler benzer g\"or\"unm\"u\c{s}t\"ur.
\end{enumerate}

\subsection{Bu Bulgular\i{}n Anlam\i{}}
\begin{itemize}
    \item Sonu\c{c} listesindeki \c{c}iftlerin b\"uy\"uk b\"ol\"um\"u ger\c{c}ek sendikasyon/yeniden yay\i{}n kaynakl\i{} duplicate kay\i{}tlard\i{}r.
    \item Google News, Yahoo/AOL/Globe gibi agregat\"or ge\c{i}\c{s}leri similarity de\u{g}erlerini y\"ukseltmektedir.
    \item ``news'' gibi tek-kelimeye d\"u\c{s}en normalize ba\c{s}l\i{}klar, duplicate tespit kalitesini olumsuz etkileyen ana risklerden biridir.
\end{itemize}

\section{Zipf Yasas\i{} Analizi}
Notebook i\c{c}inde rank-frequency da\u{g}\i{}l\i{}m\i{} log-log d\"uzleminde incelenmi\c{s}tir. Hesaplanan e\u{g}imler:
\begin{itemize}
    \item \.{I}lk 300 rank i\c{c}in e\u{g}im: \textbf{-0.628}
    \item \.{I}lk 100 rank i\c{c}in e\u{g}im: \textbf{-0.522}
\end{itemize}

Klasik Zipf davran\i{}\c{s}\i{}nda e\u{g}imin -1 civar\i{}nda olmas\i{} beklenir. Bu projede elde edilen de\u{g}erlerin daha d\"u\c{s}\"uk mutlak de\u{g}erde kalmas\i{}n\i{}n ana sebepleri:
\begin{itemize}
    \item Veri tipinin haber ba\c{s}l\i{}\u{g}\i{} olmas\i{} (k\i{}sa metinler),
    \item Veri havuzunun konu odakl\i{} (tech/business/world) toplanmas\i{},
    \item Kaynak heterojenli\u{g}i ve ba\c{s}l\i{}k stil farklar\i{}.
\end{itemize}

Bununla birlikte ilk denemelere g\"ore g\"ur\"ult\"u temizli\u{g}i sonras\i{}nda da\u{g}\i{}l\i{}m daha anlaml\i{} hale gelmi\c{s}tir (\"orne\u{g}in \texttt{new}, saf y\i{}l/say\i{} tokenlar\i{} gibi kelimeler filtrelenmi\c{s}tir).

\section{Normalizasyon \.Iyile\c{s}tirmeleri}
\texttt{analysis\_pipeline.ipynb} i\c{c}indeki normalizasyon ad\i{}m\i{}nda:
\begin{itemize}
    \item Stopword temizli\u{g}i (\texttt{nltk}),
    \item Noise kelime listesi geni\c{s}letme (\texttt{breaking, report, latest, new}),
    \item Saf say\i{}sal tokenlar\i{}n elenmesi,
    \item 3 karakterden k\i{}sa tokenlar\i{}n elenmesi
\end{itemize}
uygulanm\i{}\c{s}t\i{}r.

\section{Sonu\c{c}}
Sistem; veri toplama, birle\c{s}tirme, normalizasyon, duplicate tespiti ve Zipf temelli kalite incelemesini tek bir i\c{s} ak\i{}\c{s}\i{}nda tamamlamaktad\i{}r. Elde edilen sonu\c{c}lar pratik bir haber ba\c{s}l\i{}\u{g}\i{} veri seti i\c{c}in tutarl\i{}d\i{}r. Zipf e\u{g}imi klasik uzun edebi metinlere g\"ore d\"u\c{s}\"uk olsa da, veri tipinin do\u{g}as\i{} gere\u{g}i kabul edilebilir bir seviyededir ve temizleme ad\i{}mlar\i{}yla anlaml\i{}l\i{}k belirgin \c{s}ekilde artm\i{}\c{s}t\i{}r.

\end{document}
